\section{Construcción matemática de los conjuntos experimentales}\label{sec:conjuntos}

Los $28$ escenarios evaluados corresponden a configuraciones experimentales derivadas sistemáticamente del mismo panel de $17$ centrales mediante combinaciones de horizonte, estrategia, disponibilidad meteorológica y régimen de generalización (Material Suplementario, Sección~S3, Tabla~S4), por lo que no constituyen $28$ conjuntos de datos independientes.

\subsection{Panel base, unidades físicas y filtración causal}

Sea $\mathcal{P} = \{1,\dots,M\}$ ($M=17$) el conjunto de centrales, cada una en una zona espacial fija $g(i) \in \{1,\dots,G\}$ ($G=7$), y sea $\mathcal{T}_i$ el conjunto de días válidos de la planta $i$ entre el 1 de enero de 2024 y el 31 de octubre de 2026. Con energía diaria $y_{i,s} > 0$ ($\text{kWh/día}$) y capacidad efectiva neta $C_{i,s} > 0$ ($\text{MW}$), la cota energética diaria bajo operación nominal continua de $24\text{ h}$ es:
\begin{equation}
E^{\max}_{i,s} := 24\,000\,C_{i,s} \quad (\text{kWh/día}).
\label{eq:emax}
\end{equation}
Al cierre operativo de un origen $t \in \mathcal{T}_i$, el denominador causal vigente es $E^{\max}_{i,t}$, evitando anticipar ampliaciones posteriores $C_{i,t+h}$ ($h \ge 1$). La filtración causal $\mathcal{F}_{i,t} := \sigma(\{y_{i,s}, C_{i,s} : s \le t\},\; \{\mathbf{z}_{i,u \mid s} : s \le t,\; u > s\},\; \text{calendario})$ agrupa observaciones pasadas y pronósticos emitidos en $s \le t$ para $u > s$. En todo el manuscrito, «causal» designa estrictamente no anticipación temporal respecto a $\mathcal{F}_{i,t}$.

\subsection{Definición de objetivos y representación de entrada}

Para planta $i$, origen $t$ y horizonte $H \in \{1,7,14,21,28\}$ días, se distinguen dos objetivos:
\begin{enumerate}
    \item \textbf{Objetivo diario ($D+1$ con $H=1$ y recursivo con $H \in \{7,14,21,28\}$):} para cada paso $h \in \{1,\dots,H\}$, la respuesta física es $y_{i,t+h}$ ($H_{\mathrm{dir}} = 1$) y su factor de capacidad causal truncado es:
    \begin{equation}
    c_{i,t+h \mid t} := \clip\!\left(\frac{y_{i,t+h}}{24\,000\,C_{i,t}},\, 0,\, 1\right),
    \qquad \clip(u,0,1) := \min\{\max\{u,0\},1\}.
    \label{eq:cf_diario}
    \end{equation}
    \item \textbf{Objetivo acumulado (esquema directo con $H \in \{7,14,21,28\}$):} predice en un paso la energía acumulada sobre ventanas válidas $[t+1, t+H]$ ($H_{\mathrm{dir}} = H$):
    \begin{equation}
    Y_{i,t}^{(H)} := \sum_{h=1}^{H} y_{i,t+h},
    \qquad
    c_{i,t}^{(H)} := \clip\!\left(\frac{Y_{i,t}^{(H)}}{24\,000\,H\,C_{i,t}},\, 0,\, 1\right),
    \label{eq:y_directo}
    \end{equation}
    dividiendo el error acumulado en $\text{kWh}$ por $H$ para expresarlo en $\text{kWh/día}$.
\end{enumerate}

A partir de las medias móviles causales de $30$ y $7$ días y del último registro válido hasta $t$ ($\bar{y}^{(30)}_{i,t}$, $\bar{y}^{(7)}_{i,t}$, $y^{\mathrm{ult}}_{i,t}$), se calculan los niveles recientes acotados:
\begin{equation}
\mathrm{cf}_{30}(i,t) := \clip\!\left(\frac{\bar{y}^{(30)}_{i,t}}{E^{\max}_{i,t}}, 0, 1\right),\quad
\mathrm{cf}_{7}(i,t) := \clip\!\left(\frac{\bar{y}^{(7)}_{i,t}}{E^{\max}_{i,t}}, 0, 1\right),\quad
\mathrm{cf}_{1}(i,t) := \clip\!\left(\frac{y^{\mathrm{ult}}_{i,t}}{E^{\max}_{i,t}}, 0, 1\right),
\label{eq:cf_hist}
\end{equation}
con respaldos jerárquicos causales en los primeros días sin ventana completa (Material Suplementario, Sección~S3). Para la fecha objetivo $t+h$, el vector de predictores $\mathbf{x}_{i,t+h \mid t} \in \mathbb{R}^p$ incluye:
\begin{itemize}
    \item \textbf{Sin meteorología ($p = 7$):} cuatro armónicos anuales evaluados en $t+h$ ($\sin/\cos$ de mes y día del año) y tres términos de dinámica autorregresiva con amortiguamiento geométrico $\gamma_h := 0{,}85^{h-1}$:
    \begin{equation}
    x_5 = \mathrm{cf}_{30}(i,t),\quad
    x_6 = \gamma_h \big(\mathrm{cf}_{7}(i,t) - \mathrm{cf}_{30}(i,t)\big),\quad
    x_7 = \gamma_h \big(\mathrm{cf}_{1}(i,t) - \mathrm{cf}_{7}(i,t)\big).
    \label{eq:features_din}
    \end{equation}
    \item \textbf{Con meteorología ($p = 17$):} añade las cinco variables pronosticadas $\mathbf{z}_{i,t+h \mid t}$ (irradiación horizontal, temperatura a $2\text{ m}$, nubosidad total, humedad relativa y viento a $10\text{ m}$) y cinco anomalías causales $\mathbf{a}_{i,t+h \mid t} = \mathbf{z}_{i,t+h \mid t} - \bar{\mathbf{z}}^{(30)}_{i,t}$ respecto a la media móvil de $30$ días de los pronósticos emitidos hasta $t$.
\end{itemize}

\subsection{Particiones externas, validación interna y actualización recursiva}

Las $28$ configuraciones cruzan cuatro familias de tarea ($D+1$ con y sin clima en tres bloques trimestrales nov-2025 a jul-2026; recursivo sin clima; directo acumulado sin clima; y recursivo con GEFSv12, con orígenes no solapados tras $90$ días de calentamiento para $H \in \{7,14,21,28\}$) con dos regímenes de generalización:
\begin{itemize}
    \item \textbf{Régimen intraplanta (IP):} al evaluar la fecha u origen $t_0$ de la planta $i$, $\mathcal{D}_{\mathrm{train}}(i,t_0)$ usa únicamente observaciones de la misma planta $i$ cuyo periodo objetivo concluye en $t_0$ o antes.
    \item \textbf{Régimen espaciotemporal (ET):} validación \textit{leave-one-group-out} sobre los $G=7$ conglomerados espaciales, entrenando con $\bigcup_{j : g(j) \ne g(i)} \mathcal{F}_{j,t_0}$. Este régimen \emph{no equivale a un escenario \textit{zero-shot} sin historial local}, pues la planta evaluada $i$ sí aporta su historial causal previo hasta $t_0$ ($\mathrm{cf}_{30}$, $\mathrm{cf}_7$, $\mathrm{cf}_1$ y $\bar{\mathbf{z}}^{(30)}_{i,t_0}$) para construir sus predictores al origen.
\end{itemize}

En \texttt{SM-CFHR-HPO}, la búsqueda interna ocurre dentro de $\mathcal{D}_{\mathrm{train}}$ mediante validación cruzada causal de $3$ particiones temporales (embargo de $30$ días en $D+1$; $t_{\mathrm{end}} \le t_{\mathrm{val}}$ en multihorizonte; exclusión espacial en ET). Para horizontes recursivos ($H \in \{7,14,21,28\}$), tanto la validación interna como la evaluación externa fijan $C_{i,t_0}$ y $\bar{\mathbf{z}}^{(30)}_{i,t_0}$ al origen e iteran paso a paso ($h=1,\dots,H$) desplazando el búfer histórico de $30$ días exclusivamente con las predicciones propias $\hat{y}_{i,t_0+h \mid t_0}$ (Material Suplementario, Sección~S3, Algoritmo~S1).
